% Chapter 4: Learners and Teachers
% Rewritten against the design-lab Lean 4 kernel (FLT_Proofs.Learner.*) and the
% transformers library (TLT_Proofs): typed learner abstraction, the Gibbs variational
% principle in full, and a modern attention-learner instance of MeasurableBatchLearner.
%   14 x Profile A (vocabulary: learner types and teacher types - TABLE, not definitions)
%   bayesian_learner = Profile E (Obstruction: posterior ≠ single hypothesis)
%   gibbs_posterior = Profile C (Revelation: the Bayesian-PAC bridge)
%   meta_learner = Profile F (Open Frontier: learning to learn)
%   team_learner = Profile C (Revelation: teams > individuals)
%   optimal_teacher = Profile G (Framework Bridge → teaching_dimension)
%   minimally_adequate_teacher = Profile G (Framework Bridge → L*)

\chapter{Learners and Teachers}
\label{ch:agents}

Chapters~\ref{ch:objects} and~\ref{ch:data} established what is learned (concept classes) and how data arrives (i.i.d.\ samples, texts, queries, streams). This chapter introduces the agents that do the learning and, in some settings, the agents that provide the data.

Much of the vocabulary here is taxonomic: a learner is a function from data to hypotheses, and the field has named a dozen variants distinguished by which constraints they satisfy. These are best presented as a table. The chapter's mathematical content lies in three places: the type mismatch between Bayesian posteriors and PAC hypotheses, resolved by the Gibbs posterior, whose optimality we prove; the power of team learning; and the teacher models that connect to the teaching dimension and to Angluin's $L^\ast$. We close the opening section with a modern instance, a finite-head attention learner, that satisfies the measurability condition the PAC theory requires.

% ================================================================
% SECTION 1: The Learner Abstraction
% ================================================================

\section{The Learner Abstraction}
\label{sec:learner}

\begin{defn}[Learner]
\label{def:learner}
A \emph{learner} is a function
\[
L \colon \bigcup_{m=0}^{\infty} (X \times Y)^m \;\to\; \mathcal{H}
\]
mapping a finite data sequence to a hypothesis: having observed $S_t = ((x_1, y_1), \ldots, (x_t, y_t))$, it outputs $h_t = L(S_t) \in \mathcal{H}$. The verified development splits this single abstraction into the three paradigm-specific carriers: a \emph{batch learner} maps a whole sample to a hypothesis,\formal{FLT_Proofs.Learner.Core.BatchLearner} an \emph{online learner} carries a state updated example by example,\formal{FLT_Proofs.Learner.Core.OnlineLearner} and a \emph{Gold learner} conjectures from the data seen so far.\formal{FLT_Proofs.Learner.Core.GoldLearner}
\end{defn}

This definition is minimal. It says nothing about how $h_t$ is selected: whether the learner stores the full history or only its previous guess, whether it uses randomness, whether it may choose which $x$ to query. Each constraint yields a named variant (Figure~\ref{fig:learner-taxonomy} and Table~\ref{tab:learner-types}). One further condition is invisible at the level of types but indispensable for the statistical theory.

\begin{defn}[Measurable batch learner]
\label{def:measurable-learner}
A batch learner $L$ over a measurable domain is a \emph{measurable batch learner} if its evaluation map is jointly measurable: for every sample size $m$, the map $(S, x) \mapsto L(S)(x)$ from $(X \times Y)^m \times X$ to $Y$ is measurable.\formal{FLT_Proofs.Learner.Core.MeasurableBatchLearner}
\end{defn}

This is the minimal regularity that makes the PAC success event $\{S : \Pr_x[L(S)(x) \neq c(x)] \leq \varepsilon\}$ a measurable set, so that ``with probability $1 - \delta$'' is well-defined; without it the PAC criterion of \Cref{ch:pac} cannot even be stated~\cite{KrappWirth2024}. Every theorem in the statistical chapters carries this hypothesis silently, and the learners that matter in practice are exactly the ones for which measurability is not obvious.

\subsection{A modern instance: the attention learner}
\label{sec:attention-learner}

A transformer's attention layer is a learner in the sense of \cref{def:learner}: its parameters are fit to data, and at inference it routes each input through one of several ``heads.'' Whether this learner is admissible for the statistical theory is exactly the question of \cref{def:measurable-learner}, and the answer is yes.

\begin{prop}[A finite-head attention learner is a measurable batch learner]
\label{prop:attention-mbl}
Fix $k \geq 1$. A \emph{$k$-head attention learner} is given by a standard-Borel parameter space $\Theta$, measurable score functions $s_\theta : X \to \R^k$, measurable experts $e_\theta : X \to Y^k$, and a measurable parameter-selection rule $\theta(\cdot)$ from samples; its prediction routes $x$ to the expert of highest score, $L(S)(x) = e_{\theta(S)}\bigl(x\bigr)_{\arg\max_i s_{\theta(S)}(x)_i}$. Then $L$ is a measurable batch learner, so every PAC theorem of \Cref{ch:pac} applies to it.\formal{TLT_Proofs.Learner.MeasurableParametricAttentionLearner.ParametricFiniteHeadAttentionLearner.instMBL'}
\end{prop}

\begin{proof}
The argmax router partitions the input space into the cells $\{x : i = \arg\max_j s_\theta(x)_j\}$, each a finite intersection of the measurable score-comparison sets $\{x : s_\theta(x)_i \geq s_\theta(x)_j\}$; hence the routing map $(\theta, x) \mapsto \mathrm{route}_\theta(x)$ is measurable.\formal{TLT_Proofs.Routing.MeasurableScoreRouting.FiniteScoreRouterCode.route_measurable} Composing the measurable router with the measurable experts gives a jointly measurable evaluation $(\theta, x) \mapsto e_\theta(x)_{\mathrm{route}_\theta(x)}$. Precomposing with the measurable selection $S \mapsto \theta(S)$ in the first coordinate preserves joint measurability, so $(S,x) \mapsto L(S)(x)$ is measurable, which is exactly \cref{def:measurable-learner}. The whole construction is a parametric learner family, and every parametric family with a measurable evaluation and selection is a measurable batch learner.\formal{TLT_Proofs.Learner.MeasurableParametricAttentionLearner.ParametricLearnerFamily.instMeasurableBatchLearner}
\end{proof}

The same holds for the softmax-top-1 and multi-head argmax routers, and for the scaled dot-product attention of a full transformer.\formal{TLT_Proofs.Routing.ScaledDotProductScoreRouter.attentionRouting_wellBehaved} The abstract learner of \cref{def:learner} therefore covers the dominant modern architectures, and places them in the measurable sub-class on which the PAC apparatus operates. We return to the measurability boundary, where attention routing \emph{fails} to be Borel, in \Cref{ch:separations}.

% ================================================================
% SECTION 2: Taxonomy Table
% ================================================================

\section{Taxonomy of Learner Types}
\label{sec:learner-taxonomy}

\begin{figure}[ht]
\centering
\begin{tikzpicture}[
    font=\small,
    root/.style={draw, very thick, rounded corners, fill=defblue!15, font=\small\bfseries,
                 minimum width=2.4cm, minimum height=0.6cm, align=center},
    pred/.style={draw, semithick, rounded corners, fill=layergray, font=\small,
                 minimum width=2.5cm, minimum height=0.7cm, align=center},
    sink/.style={draw, thick, rounded corners, align=center, font=\scriptsize,
                 text width=3.3cm, inner sep=4pt},
    gedge/.style={-{Stealth[length=2mm]}, sepgreen!60!black, semithick},
    oedge/.style={-{Stealth[length=2mm]}, edgeorange!80!black, semithick, dashed},
    meet/.style={notegray, dotted, thick},
    elbl/.style={font=\scriptsize\itshape, fill=white, inner sep=1pt, text=black!65}
]
% --- root anchor (held-constant Will) ---
\node[root] (root) at (0,3.0) {Learner\\[-1pt]{\scriptsize\mdseries one learner, many constraints}};
% --- predicate spine (verified typed predicates, solid-outlined nodes) ---
\node[pred] (cons)  at (0,1.55) {Consistent\\[-2pt]{\tiny\ttfamily IsConsistent}};
\node[pred] (consv) at (0,0.45) {Conservative\\[-2pt]{\tiny\ttfamily IsConservative}};
\node[pred] (sd)    at (0,-0.65) {Set-driven\\[-2pt]{\tiny\ttfamily IsSetDriven}};
\node[pred] (it)    at (0,-1.75) {Iterative\\[-2pt]{\tiny\ttfamily IsIterative}};
% --- faint lattice-meet connectors (constraints on one learner) ---
\draw[meet] (root) -- (cons);
\draw[meet] (cons) -- (consv);
\draw[meet] (consv) -- (sd);
\draw[meet] (sd) -- (it);
% --- regime sinks ---
\node[sink, sepgreen!60!black, fill=sepgreen!8] (pac) at (-4.7,0.1)
  {\textbf{In PAC: constraints are FREE}\\[2pt]learnable by a consistent proper ERM
   (\leanname{ermLearn}); i.i.d.\ makes order irrelevant\\[2pt]
   verified: \leanname{erm_consistent_realizable}};
\node[sink, edgeorange!80!black, fill=edgeorange!8, dashed] (gold) at (4.7,-0.55)
  {\textbf{In Gold/online: constraints BITE}\\[2pt]conservative from informant, not text;
   set-driven forces unbounded memory on iterative\\[2pt]
   literature: Gold; Angluin; K\"otzing--Palenta};
% --- green solid edges: free in PAC ---
\draw[gedge] (cons) -- (pac);
\draw[gedge] (sd)   -- node[elbl, above, pos=0.45]{i.i.d.} (pac);
\draw[gedge] (it)   -- (pac);
% --- orange dashed edges: biting in Gold ---
\draw[oedge] (consv) -- (gold);
\draw[oedge] (sd)    -- node[elbl, above, pos=0.45]{memory} (gold);
\draw[oedge] (it)    -- (gold);
% --- legend ---
\node[font=\scriptsize, text=black!60, align=center] at (0,-3.05)
  {Solid green: free in PAC (kernel-verified).\quad Dashed orange: biting in Gold (literature).\quad
   Boxes: verified predicates.};
\end{tikzpicture}
\caption[Lattice of learner constraints: free in PAC, biting in Gold]{The learner taxonomy as a
lattice of typed constraints on the base learner of \cref{def:learner}. The same four predicates read
two ways: in PAC they are free, since every PAC-learnable class is learnable by a consistent, proper
empirical-risk minimizer (\leanname{erm_consistent_realizable}) and the i.i.d.\ assumption makes order
irrelevant; in Gold's framework they bite, separating conservative from text-only learners and forcing
unbounded memory on iterative ones. The named data-access and randomness variants are in
\cref{tab:learner-types}.}
\label{fig:learner-taxonomy}
\end{figure}

The named variants are constraints on \cref{def:learner}, and each is a typed predicate on the learner in the verified development: consistency,\formal{FLT_Proofs.Learner.Properties.IsConsistent} conservativeness,\formal{FLT_Proofs.Learner.Properties.IsConservative} set-drivenness,\formal{FLT_Proofs.Learner.Properties.IsSetDriven} and iterativeness.\formal{FLT_Proofs.Learner.Properties.IsIterative} Table~\ref{tab:learner-types} gives the precise constraints.

\begin{table}[ht]
\centering
\small
\caption{Named learner types. Each row gives the defining constraint and the chapter where the concept plays its primary role. These are not independent definitions; a single learner may satisfy several constraints simultaneously (e.g., a conservative, set-driven, passive learner).}
\label{tab:learner-types}
\begin{tabularx}{\textwidth}{>{\bfseries\small}l X l}
\toprule
\textbf{Name} & \textbf{Defining Constraint} & \textbf{Primary Role} \\
\midrule
Passive &
  Receives data without choosing which instances to observe. The standard assumption in PAC and Gold settings. &
  Ch.~\ref{ch:pac}, \ref{ch:gold} \\[4pt]

Active &
  Chooses which instances to query. Performance measured by \emph{label complexity}: the number of labels requested, not the number of unlabeled instances seen. &
  Ch.~\ref{ch:exact} \\[4pt]

Conservative &
  Changes hypothesis only when forced by inconsistency: if $h_t$ is consistent with $S_{t+1}$, then $h_{t+1} = h_t$. &
  Ch.~\ref{ch:gold} \\[4pt]

Consistent &
  Always outputs a hypothesis consistent with all data seen so far: $h_t(x_i) = y_i$ for all $i \leq t$. &
  Ch.~\ref{ch:pac}, \ref{ch:gold} \\[4pt]

Set-driven &
  Output depends only on the \emph{set} $\{(x_i, y_i)\}$ of data, not on the order of presentation. Formally, $L(S) = L(S')$ whenever $S$ and $S'$ are permutations of each other. &
  Ch.~\ref{ch:gold} \\[4pt]

Iterative &
  Hypothesis at step $t$ depends only on $h_{t-1}$ and the new datum $(x_t, y_t)$, not on the full history $S_t$. Memory is bounded: the learner's state is a single hypothesis. &
  Ch.~\ref{ch:gold}, \ref{ch:online} \\[4pt]

Probabilistic &
  May use internal randomness in selecting hypotheses. The success criterion (PAC, identification, mistake bound) is then required to hold with high probability over both the data and the learner's coin flips. &
  Ch.~\ref{ch:pac} \\[4pt]

With advice &
  Receives $b$ bits of advice in addition to data. The advice may encode information about the target class or the distribution, and $b$ parameterizes the amount of side information. &
  Ch.~\ref{ch:exact} \\[4pt]

\bottomrule
\end{tabularx}
\end{table}

These distinctions generate strict separations in learning power within Gold's framework: there are classes identifiable by a conservative learner from informant but not from text, and classes identifiable by a set-driven learner that require unbounded memory for an iterative learner (\Cref{ch:gold,ch:separations}).

In the PAC setting many of the distinctions collapse, and the collapse has the status of a proved theorem: every PAC-learnable class is learnable by a consistent, proper learner (the empirical-risk minimizer, which returns any hypothesis consistent with the sample and succeeds once uniform convergence holds (\cref{thm:uc-pac})).\formal{FLT_Proofs.Complexity.Generalization.ermLearn} Because the i.i.d.\ assumption makes example order irrelevant, the set-driven/order-sensitive distinction also vanishes. The taxonomy has its sharpest teeth in the Gold and online paradigms, where the learner's memory and update policy genuinely affect what is learnable.

% ================================================================
% SECTION 3: The Bayesian–PAC Bridge
% ================================================================

\section{The Bayesian--PAC Bridge}
\label{sec:bayesian-pac}

The learner types in Table~\ref{tab:learner-types} all output a single hypothesis $h \in \mathcal{H}$. The Bayesian learner does not.

\begin{defn}[Bayesian Learner]
\label{def:bayesian-learner}
A \emph{Bayesian learner} maintains a prior $\pi$ over $\mathcal{H}$ and, given data $S$, outputs the posterior
\[
Q(h \mid S) \;=\; \frac{\pi(h) \cdot P(S \mid h)}{\sum_{h' \in \mathcal{H}} \pi(h') \cdot P(S \mid h')},
\]
a whole distribution over hypotheses, with every $h$ retained at its posterior weight.\formal{FLT_Proofs.Learner.Bayesian.BayesianLearner}
\end{defn}

\subsection{The type mismatch}

PAC learning requires a learner that outputs a single $h \in \mathcal{H}$ with $\risk_D(h) \leq \risk_D(h^\ast) + \varepsilon$ at confidence $1 - \delta$. The Bayesian learner outputs a probability measure on $\mathcal{H}$. These are different types:
\[
\text{PAC learner:} \quad (X \times Y)^m \to \mathcal{H},
\qquad
\text{Bayesian learner:} \quad (X \times Y)^m \to \Delta(\mathcal{H}).
\]
The naive resolution, outputting the MAP hypothesis $\hat{h} = \arg\max_h Q(h \mid S)$, discards most of the posterior and yields suboptimal generalization. The verified development records both the MAP read-off and the conversion of a Bayesian learner to an ordinary batch learner.\formal{FLT_Proofs.Bridge.bayesianToBatch} The sharper resolution is the Gibbs posterior.

\begin{obstbox}
\textbf{Bayesian $\to$ PAC type mismatch.} The Bayesian learner's output type ($\Delta(\mathcal{H})$) does not match the PAC learner's output type ($\mathcal{H}$). The PAC-Bayes theorem (\Cref{ch:generalization}) bounds the risk of a \emph{distribution} over hypotheses; extracting a single hypothesis with the bound intact requires an explicit randomized construction: the Gibbs posterior.
\end{obstbox}

\subsection{The Gibbs posterior and its optimality}

\begin{defn}[Gibbs Posterior]
\label{def:gibbs-posterior}
Given a prior $\pi$ over $\mathcal{H}$, the empirical risk $\emprisk_S$, and an inverse temperature $\lambda > 0$, the \emph{Gibbs posterior} is
\[
Q_\lambda(h) \;\propto\; \pi(h) \cdot \exp\!\bigl(-\lambda \,\emprisk_S(h)\bigr).
\]
The \emph{Gibbs classifier} draws $h \sim Q_\lambda$ and predicts with $h$; its expected risk is $\E_{h \sim Q_\lambda}[\risk_D(h)]$.\formal{FLT_Proofs.Learner.Bayesian.GibbsPosterior}
\end{defn}

The PAC-Bayes theorem~\cite{McAllester1999} states that for any prior $\pi$ and any posterior $Q$, with probability $1 - \delta$ over $S \sim D^m$,
\[
\E_{h \sim Q}[\risk_D(h)] \;\leq\; \E_{h \sim Q}[\emprisk_S(h)] + \sqrt{\frac{\KL(Q \,\|\, \pi) + \ln(m/\delta)}{2m}},
\]
proved for finite hypothesis spaces in the verified development.\formal{FLT_Proofs.Theorem.PACBayes.pac_bayes_finite} The bound holds for \emph{every} $Q$, and so leaves open which $Q$ the learner should use. The Gibbs posterior is the answer, by a variational argument.

\begin{prop}[The Gibbs posterior minimizes the PAC-Bayes objective {\cite{Catoni2007}}]
\label{prop:gibbs-variational}
Let $\mathcal{H}$ be finite and $\lambda > 0$. Among all distributions $Q$ on $\mathcal{H}$, the objective
\[
F_\lambda(Q) \;=\; \lambda\,\E_{h \sim Q}[\emprisk_S(h)] + \KL(Q \,\|\, \pi)
\]
is minimized uniquely by the Gibbs posterior $Q_\lambda$, with minimum value $-\ln Z_\lambda$ where $Z_\lambda = \sum_{h} \pi(h)\,e^{-\lambda \emprisk_S(h)}$.
\end{prop}

\begin{proof}
Write $Q_\lambda(h) = \pi(h)\,e^{-\lambda \emprisk_S(h)} / Z_\lambda$. For any distribution $Q$ on $\mathcal{H}$,
\[
\KL(Q \,\|\, \pi) + \lambda\,\E_{h\sim Q}[\emprisk_S(h)]
= \sum_h Q(h)\Bigl[\ln \tfrac{Q(h)}{\pi(h)} + \lambda\,\emprisk_S(h)\Bigr]
= \sum_h Q(h)\,\ln \frac{Q(h)}{\pi(h)\,e^{-\lambda \emprisk_S(h)}}.
\]
Substituting $\pi(h)\,e^{-\lambda \emprisk_S(h)} = Z_\lambda\,Q_\lambda(h)$ and splitting the logarithm,
\[
= \sum_h Q(h)\,\ln \frac{Q(h)}{Z_\lambda\,Q_\lambda(h)}
= \sum_h Q(h)\,\ln \frac{Q(h)}{Q_\lambda(h)} - \ln Z_\lambda
= \KL(Q \,\|\, Q_\lambda) - \ln Z_\lambda .
\]
By Gibbs' inequality $\KL(Q \,\|\, Q_\lambda) \geq 0$, with equality if and only if $Q = Q_\lambda$. Hence $F_\lambda(Q) \geq -\ln Z_\lambda$, with equality exactly at $Q = Q_\lambda$.
\end{proof}

The PAC-Bayes bound's right-hand side is, up to the monotone square root and constants, an instance of $F_\lambda$: minimizing the complexity-penalized empirical risk over posteriors yields $Q_\lambda$. The bound thus names the exact posterior to use, which is the precise sense in which the Gibbs posterior converts a Bayesian learner into a PAC learner.\formal{FLT_Proofs.Theorem.PACBayes.gibbs_bound_of_pointwise}

\begin{remark}[Gibbs $\neq$ MAP]
\label{rmk:gibbs-vs-map}
The Gibbs posterior draws a \emph{random} hypothesis weighted by empirical performance; MAP takes the \emph{mode}. In high-dimensional settings the Gibbs posterior gives tighter guarantees because PAC-Bayes controls expected risk under the full posterior, while a MAP bound requires uniform convergence over a point estimate. As $\lambda \to \infty$ the Gibbs posterior concentrates on the empirical risk minimizers and recovers MAP in the limit; for finite $\lambda$ it strictly hedges, which is the source of the improved bound.
\end{remark}

The full development of PAC-Bayes bounds, including the McAllester and Catoni variants, appears in \Cref{ch:generalization}.

% ================================================================
% SECTION 4: Teams and Meta-Learning
% ================================================================

\section{Teams and Meta-Learning}
\label{sec:teams-meta}

\subsection{Team learning}

The team model asks what happens when several learners collaborate, where success requires only that \emph{at least one} member converges to a correct hypothesis.

\begin{defn}[Team Learner]
\label{def:team-learner}
A \emph{team} of size $k$ is a tuple of learners $(L_1, \ldots, L_k)$. The team \emph{identifies} $c \in \mathcal{C}$ if at least one $L_i$ converges to a correct hypothesis; $\mathcal{C}$ is \emph{team-identifiable} if some finite team identifies every $c \in \mathcal{C}$.\formal{FLT_Proofs.Learner.Properties.TeamLearner}
\end{defn}

\begin{thm}[Teams $>$ Individuals {\cite{Smith1982}}]
\label{thm:team-power}
There exist concept classes identifiable by a team of size~2 but by no single learner.
\end{thm}

The team members receive the \emph{same} data; they neither partition the instance space nor communicate, so the power comes from the disjunctive success criterion alone: ``at least one succeeds'' is a weaker demand than ``the one learner succeeds,'' and the weaker demand enlarges the learnable classes. The separation lives inside Gold's standard framework, with ordinary recursively enumerable languages as the witnesses. The proof builds a class on which any single learner is forced into infinitely many mind changes by diagonalization (as in the text/informant separation of \Cref{ch:data}), while two learners with complementary strategies (one guessing large languages, one guessing small) guarantee that one converges. The argument appears in \Cref{ch:separations}.

\begin{remark}[Teams and probabilistic learners]
\label{rmk:team-prob}
A team of deterministic learners is strictly stronger than a single deterministic learner (\cref{thm:team-power}). A probabilistic learner succeeding with probability $> 1/2$ can simulate a $k$-team by running $k$ internal copies; but in the Gold setting, where success means convergence with certainty, this simulation is unavailable, and the team hierarchy is strict.
\end{remark}

\subsection{Meta-learning}

\begin{defn}[Meta-Learner]
\label{def:meta-learner}
A \emph{meta-learner} operates over a family of tasks $\{\mathcal{C}_j\}$ drawn from an environment $\mathcal{E}$: from experience on $\mathcal{C}_1, \ldots, \mathcal{C}_{n-1}$ it produces a bias (a hypothesis space or prior) that performs well on a fresh task $\mathcal{C}_n$ from the same environment.\formal{FLT_Proofs.Theorem.Extended.MetaLearnerPAC}
\end{defn}

Baxter~\cite{Baxter2000} established the first PAC-style bounds for meta-learning: if $\mathcal{E}$ has bounded complexity (covering numbers over a family of hypothesis spaces) and the meta-learner sees $n$ tasks of $m$ examples each, the expected excess risk on a new task is
\[
O\!\left(\sqrt{\frac{\mathrm{complexity}(\mathcal{E})}{n}} + \sqrt{\frac{\VC(\mathcal{H})}{m}}\right),
\]
a \emph{task-level} term (estimating the environment from $n$ tasks) plus a \emph{within-task} term (learning one task from $m$ examples). This two-level decomposition is the defining feature of meta-learning bounds, and whether a single combinatorial dimension characterizes meta-learnability is open (\cref{op:meta-char}).

% ================================================================
% SECTION 5: Teachers and the Teaching Game
% ================================================================

\section{Teachers and the Teaching Game}
\label{sec:teachers}

In the models above, data arrives passively or through the learner's queries. The \emph{teacher} reverses the direction: an agent chooses examples to help (or hinder) the learner.

\begin{table}[ht]
\centering
\small
\caption{Teacher types. Each row gives the teacher's strategy, the induced data model, and the chapter where the concept plays its primary role.}
\label{tab:teacher-types}
\begin{tabularx}{\textwidth}{>{\bfseries\small}l X l}
\toprule
\textbf{Name} & \textbf{Strategy and Key Property} & \textbf{Primary Role} \\
\midrule
Random &
  The \emph{random teacher}\index{random teacher} draws examples i.i.d.\ from a distribution $D$. The standard PAC assumption; the teacher has no intention, and ``teaching'' is an anthropomorphic misnomer for sampling. &
  Ch.~\ref{ch:pac} \\[4pt]

Adversarial &
  Chooses worst-case examples adaptively. The online model treats the data stream as an adversarial teacher; performance is worst-case mistake count. &
  Ch.~\ref{ch:online} \\[4pt]

Optimal &
  Selects the smallest set of examples from which the target is uniquely identified. The size of this set is the \emph{teaching dimension} of the concept within the class. &
  Ch.~13 \\[4pt]

Minimally adequate &
  Angluin's model: answers membership and equivalence queries. ``Minimally adequate'' because MQ + EQ is the minimal oracle configuration that makes DFAs efficiently learnable. &
  Ch.~\ref{ch:exact} \\[4pt]
\bottomrule
\end{tabularx}
\end{table}

\subsection{The optimal teacher and teaching dimension}

\begin{defn}[Teaching Dimension]
\label{def:teaching-dim}
The \emph{teaching dimension} of $c \in \mathcal{C}$ is the least number of labeled examples that uniquely identify $c$ within $\mathcal{C}$:
\[
\mathrm{TD}(c, \mathcal{C}) = \min\{|S| : S \subseteq X \times Y,\; c \text{ is the only } c' \in \mathcal{C} \text{ consistent with } S\}.
\]
The teaching dimension of the class is $\mathrm{TD}(\mathcal{C}) = \max_{c \in \mathcal{C}} \mathrm{TD}(c, \mathcal{C})$, and the \emph{optimal teacher} for $c$ achieves $\mathrm{TD}(c, \mathcal{C})$.\formal{FLT_Proofs.Complexity.Structures.TeachingDim}
\end{defn}

The teaching dimension measures the cooperative communication complexity of learning: how many examples a helpful teacher must provide so that a learner who knows $\mathcal{C}$ can identify the target. It is developed in full in \Cref{ch:ordinals}, where it meets the recursive teaching dimension and the teaching--learning duality.

\subsection{The minimally adequate teacher}

\begin{defn}[Minimally Adequate Teacher {\cite{Angluin1988}}]
\label{def:mat-agents}
A \emph{minimally adequate teacher} (MAT) for $\mathcal{C}$ provides a membership oracle $\mathrm{MQ}$ and an equivalence oracle $\mathrm{EQ}$ (\Cref{ch:data}). A class is MAT-learnable if some algorithm exactly identifies any $c \in \mathcal{C}$ using polynomially many queries.
\end{defn}

Angluin's $L^\ast$ algorithm (\cref{thm:lstar}) makes DFAs MAT-learnable. The teacher is ``minimal'' precisely: dropping either oracle makes the class unlearnable in polynomial time (\cref{prop:mq-alone,prop:eq-alone}). The full treatment of $L^\ast$, including the observation table and the Myhill--Nerode connection, is in \Cref{ch:exact}.

% ================================================================
% SECTION 6: Deferred Agents
% ================================================================

\section{Deferred Agent Types}
\label{sec:deferred}

Three agent concepts in the knowledge graph, \nde{synthesizer} (a program-synthesis agent), \nde{verifier} (a correctness checker for synthesizer output), and \nde{llm\_critic} (a large language model used as a critic in a learning loop), belong to the application layer. They are built \emph{from} the learner and teacher primitives above, and their treatment is deferred to \Cref{ch:frontiers}, in the context of program synthesis, neurosymbolic learning, and LLM-in-the-loop verification.

% ================================================================
% SECTION 7: Summary
% ================================================================

\section{What This Chapter Established}
\label{sec:ch4-summary}

The taxonomy rests on three substantive points.

\begin{enumerate}[leftmargin=*,itemsep=3pt]
\item \textbf{Most learner types are constraints on one abstraction.} The eight variants in Table~\ref{tab:learner-types} are orthogonal constraints (data access, state management, hypothesis selection) on the base type $L \colon (X \times Y)^\ast \to \mathcal{H}$, each a typed predicate in the development. The measurability condition (\cref{def:measurable-learner}) is the silent hypothesis of the statistical chapters, and a finite-head attention learner satisfies it (\cref{prop:attention-mbl}), placing transformers inside the abstraction on which the PAC apparatus operates.

\item \textbf{The Bayesian--PAC bridge is a variational fact.} The Bayesian learner outputs a posterior $Q \in \Delta(\mathcal{H})$; PAC learning wants a single $h \in \mathcal{H}$. The Gibbs posterior resolves this as the unique minimizer of the complexity-penalized empirical risk (\cref{prop:gibbs-variational}), which is why it is the right posterior to feed the PAC-Bayes bound.

\item \textbf{Teams are strictly more powerful than individuals.} In the Gold setting, two learners identify classes no single learner can (\cref{thm:team-power}); the mechanism is disjunctive success alone, the same data reaching both members, who never split the instance space or exchange a message.
\end{enumerate}

Teachers range from passive (random) through adversarial to cooperative (optimal, minimally adequate), connecting this chapter to the teaching dimension of \Cref{ch:ordinals} and the exact learning algorithms of \Cref{ch:exact}.

\section{Exercises and Open Problems}
\label{sec:ch4-open}

Each problem is anchored to a concept above and tagged: a \emph{known unknown} (\textsc{ku}) is an articulated open question; an \emph{unknown known} (\textsc{uk}) is a result in the literature or the verified development not yet absorbed into a clean general statement. Verified handles are named where they exist.

\begin{openprob}[A combinatorial dimension for meta-learning, \textsc{ku}]
\label{op:meta-char}
No analogue of the fundamental theorem of PAC learning exists for meta-learning: Baxter's bounds are stated through covering numbers of the environment, and no single dimension is known to stand in their place. Find a single complexity measure of the environment $\mathcal{E}$ that characterizes meta-learnability the way $\VC$ characterizes PAC learnability, building on the typed \leanname{FLT_Proofs.Theorem.Extended.MetaLearnerPAC}, or prove no such measure exists.
\end{openprob}

\begin{openprob}[A combinatorial characterization of MAT-learnability, \textsc{ku}]
\label{op:mat-char}
Which classes are learnable from a minimally adequate teacher (\cref{def:mat-agents}) in polynomial time? Beyond DFAs the picture is partial: decision trees, certain Boolean formulas, and residual finite-state automata are known positive. Give a combinatorial characterization analogous to the VC dimension for PAC learning, or show the query model admits none.
\end{openprob}

\begin{openprob}[Capacity of attention learners, \textsc{uk}]
\label{op:attention-capacity}
\Cref{prop:attention-mbl} shows a finite-head attention learner is a measurable batch learner, so the PAC theory applies. Determine its sample complexity: bound the VC or Rademacher complexity of the $k$-head attention hypothesis class in terms of $k$, the parameter dimension, and the score temperature, using the chaining and covering-number machinery already proved for attention (\leanname{TLT_Proofs.Capacity.Chaining.RademacherSymmetrization}), and decide whether the dependence on $k$ is logarithmic or linear.
\end{openprob}

\begin{openprob}[Optimal temperature for the Gibbs posterior, \textsc{ku}]
\label{op:gibbs-temperature}
\Cref{prop:gibbs-variational} fixes the optimal posterior for each $\lambda$. Determine the $\lambda$ that minimizes the resulting PAC-Bayes risk bound as a function of sample size $m$ and prior $\pi$, and characterize when the optimal $\lambda$ is finite (strict hedging) versus $\infty$ (MAP), sharpening \cref{rmk:gibbs-vs-map}.
\end{openprob}

\begin{openprob}[The team hierarchy, \textsc{ku}]
\label{op:team-hierarchy}
\Cref{thm:team-power} separates teams of size~2 from single learners. Is the team hierarchy strict at every level: does a team of size $k+1$ identify classes that no team of size $k$ can, and is there a single parameter of $\mathcal{C}$ that determines the least team size required? Formalize the hierarchy on top of \leanname{FLT_Proofs.Learner.Properties.TeamLearner}.
\end{openprob}

\begin{openprob}[Teaching dimension versus VC dimension, \textsc{ku}]
\label{op:td-vc}
The teaching dimension (\cref{def:teaching-dim}) and the VC dimension both measure a class, but neither bounds the other in general. Characterize the classes for which $\mathrm{TD}(\mathcal{C}) = O(\VC(\mathcal{C}))$, and determine the extremal gap between the two over classes of fixed VC dimension, connecting to the recursive teaching dimension of \Cref{ch:ordinals}.
\end{openprob}

\begin{openprob}[Measurable selection for Bayesian learners, \textsc{uk}]
\label{op:bayes-measurable}
\Cref{def:measurable-learner} is the regularity the PAC theory needs, and \cref{def:bayesian-learner} outputs a posterior. Determine conditions on the prior and likelihood under which the MAP read-off (\leanname{FLT_Proofs.Bridge.bayesianToBatch}) is a \emph{measurable} batch learner, so that Bayesian learners inherit the PAC apparatus without passing through the Gibbs construction.
\end{openprob}

\begin{openprob}[Order-sensitivity outside i.i.d., \textsc{ku}]
\label{op:order-sensitivity}
In the PAC setting the set-driven/iterative distinction collapses because example order is irrelevant. Quantify the cost of order-sensitivity in the online and Gold settings: give a class on which a set-driven learner needs exponentially less memory than any iterative learner, and characterize the memory-versus-order trade-off.
\end{openprob}

\begin{openprob}[A measurability dimension for learners, \textsc{uk}]
\label{op:measurability-dimension}
\Cref{def:measurable-learner} is binary: a learner is measurable or not. The attention router is measurable (\cref{prop:attention-mbl}) but lies one quantifier away from a non-Borel construction (\Cref{ch:separations}). Define a graded measurability complexity for learner families that places attention precisely, and relate it to the analytic-but-not-Borel boundary the development witnesses for attention routing.
\end{openprob}

\begin{openprob}[Teams of attention learners, \textsc{uk}]
\label{op:team-attention}
Combine \cref{thm:team-power,prop:attention-mbl}: does a team of finite-head attention learners (a mixture-of-experts ensemble selected by the disjunctive criterion) identify classes that no single attention learner can, and is the resulting team itself a measurable batch learner? This connects the team hierarchy to the routed-expert architectures of modern practice.
\end{openprob}

\subsection*{Counterfactual exercises}

Each exercise imposes a constraint from Figure~\ref{fig:learner-taxonomy} on the base learner and asks how its effect differs between the PAC and Gold paradigms.

\begin{enumerate}[leftmargin=*,itemsep=3pt]
\item Impose consistency (\leanname{FLT_Proofs.Learner.Properties.IsConsistent}) on a PAC learner. Show it costs nothing: every PAC-learnable class is learnable by a consistent, proper empirical-risk minimizer. Trace the result through \leanname{erm_consistent_realizable}, \leanname{empError_zero_iff_consistent}, and \leanname{fundamental_theorem}.

\item Read conservativeness (\leanname{FLT_Proofs.Learner.Properties.IsConservative}) in the Gold setting. Some classes a conservative learner identifies from an informant it cannot identify from text (\cite{Angluin1980,Gold1967,CaseSmith1983}). Reconstruct the separation.

\item Take consistency alone in both paradigms. It is free in PAC (\leanname{erm_consistent_realizable}) yet in Gold can force a learner into infinitely many mind changes. Exhibit a class realizing the Gold side, and explain why the i.i.d.\ assumption removes the cost in PAC.

\item Compare set-drivenness (\leanname{FLT_Proofs.Learner.Properties.IsSetDriven}) with iterativeness (\leanname{FLT_Proofs.Learner.Properties.IsIterative}) in Gold and online learning. A set-driven class can force unbounded memory on any iterative learner (\cite{KotzingPalenta2016,Gold1967}); this separation is open in the kernel, recorded as \cref{op:order-sensitivity}. Move the same pair into PAC and show the distinction collapses because example order is irrelevant.

\item Read properness under cryptographic hardness, where the spent resource is computation time, not information. Properness preserves the finite-VC PAC class yet can cost strictly more computation (\cite{PittValiant1988}). Identify which axis the separation lives on.

\item Add internal randomness, the probabilistic variant (\leanname{FLT_Proofs.Learner.Properties.ProbabilisticLearner}). In PAC randomness derandomizes at no asymptotic cost; in the Gold limit probabilistic identification is a distinct success notion. Contrast the two.

\item Replace a single learner by a team (\leanname{FLT_Proofs.Learner.Properties.TeamLearner}). Unlike consistency, the team relaxation keeps its bite in Gold, where two learners identify classes no single learner can (\cref{thm:team-power}), and survives the move to PAC uncollapsed. Connect this to the team hierarchy of \cref{op:team-hierarchy}.

\item Ask whether ``impose consistency and properness'' is a closure operator under which PAC learnability is invariant and idempotent, so that a constraint is free exactly when it lies below the closure's fixed points. Such a closure would classify, by one Galois-style law, which constraints collapse in which paradigm. No such adjunction is defined; record what a definition would have to supply.
\end{enumerate}
